\chapter{Objetivos}

\label{chap:objetivos}

\section{Objetivo General}
El objetivo principal de este Trabajo de Fin de Grado es desarrollar una plataforma de trading algorítmico y simulación (backtesting) para el mercado de criptomonedas, capaz de integrar dinámicamente indicadores técnicos personalizados y modelos de Inteligencia Artificial para la predicción direccional del precio.

\section{Objetivos Específicos}

Para alcanzar el objetivo general propuesto, se establecen los siguientes objetivos específicos, alineados con las fases del ciclo de vida del desarrollo de software: Estos objetivos se derivan de la problemática descrita en el Capítulo \ref{chap:introduccion} y estructuran la validación técnica posterior en los Capítulos \ref{chap:metodologia}, \ref{chap:desarrollo} y \ref{chap:conclusiones}.

\begin{itemize}
    \item \textbf{Estudiar} el estado del arte del trading cuantitativo, analizando la formulación matemática de los indicadores técnicos básicos y los paradigmas actuales de \textit{Machine Learning} aplicados a series temporales financieras.
    
    \item \textbf{Analizar} los requisitos funcionales y no funcionales necesarios para la construcción de un sistema financiero de alta concurrencia.
    
    \item \textbf{Diseñar} una arquitectura de software modular y distribuida que desacople la lógica de orquestación transaccional del motor de entrenamiento e inferencia estadística.
    
    \item \textbf{Desarrollar} un motor de ingesta y sincronización de datos de mercado automatizado capaz de procesar, limpiar y almacenar masivamente velas transaccionales (\textit{candlesticks}) desde fuentes externas (APIs).
    
    \item \textbf{Implementar} un pipeline de Inteligencia Artificial paramétrico que permita al usuario configurar, entrenar y evaluar múltiples algoritmos predictivos inyectando hiperparámetros de forma dinámica.
    
    \item \textbf{Evaluar y validar} la robustez y efectividad de la plataforma mediante la ejecución de simulaciones históricas (\textit{backtesting}), analizando métricas de rendimiento computacional y financiero (tasa de acierto, factor de beneficio y reducción máxima o \textit{drawdown}).
\end{itemize}
